\apendice{Protocolo de avaliação}

% TODO: incluir a versão final do protocolo, com etapas, parâmetros e critérios
% de interpretação suficientes para reprodução ou auditoria.

\apendice{Bateria de perguntas e conjunto de referência}

% TODO: incluir ou apontar para a versão completa e congelada dos itens,
% categorias, respostas esperadas e evidências documentais.

\apendice{Rubrica de avaliação das respostas}

% TODO: incluir instruções, escala, exemplos de calibração e regras de decisão.

\apendice{Instrumentos e resultados complementares}

% TODO: incluir prompts, formulário dos avaliadores e tabelas extensas que não
% pertençam ao fluxo principal da argumentação.

\apendice{Reprodução da coleta e localização dos artefatos}
\label{ap:reproducao}

O código de coleta e suas dependências estão disponíveis no repositório do pesquisador\footnote{\url{https://github.com/fabriciosantana/mcdia/tree/fc6c67ec60ab8442db3ae9bef53898d8bf9f6350/13-dissertacao/scripts}.}. A referência fixa identifica a versão consultada nesta redação, sem pressupor que o ramo principal permanecerá inalterado. O script consulta dados públicos do Senado em janelas mensais, preserva os lotes intermediários e consolida os registros por código de pronunciamento. Para executar em ambiente Python, a partir da raiz do repositório:

\begin{verbatim}
python -m venv .venv
source .venv/bin/activate
python -m pip install -r 13-dissertacao/scripts/requirements.txt
python 13-dissertacao/scripts/01_preparar_base_discursos_batch.py \
  --data-inicio 2019-02-01 --data-fim 2023-01-31 \
  --diretorio-saida /tmp/coleta-discursos-reproducao \
  --tamanho-lote-textos 250 --trabalhadores 5 \
  --tentativas 10 --backoff 1 --pausa-entre-lotes 2
\end{verbatim}

O diretório de saída deve ser separado dos dados originais. Os lotes existentes são reutilizados na retomada; uma nova coleta independente requer um diretório vazio. O arquivo consolidado será denominado \path{discursos_2019-02-01_2023-01-31.parquet}. Para preservar os resultados fora de uma sessão temporária, deve-se substituir o caminho de exemplo por armazenamento persistente. As opções adicionais são apresentadas pelo argumento \texttt{--help}.

Reexecutar a coleta reproduz o procedimento, mas não garante identidade binária com uma coleta passada: a fonte pública e as dependências podem mudar. Para reproduzir a análise da versão congelada, deve-se obter o arquivo na revisão \texttt{v1.1.1}\footnote{\url{https://huggingface.co/datasets/fabriciosantana/discursos-senado-legislatura-56/tree/v1.1.1/data/full}.} e conferir seu SHA-256:
\begin{verbatim}
sha256sum discursos_2019-02-01_2023-01-31.parquet
\end{verbatim}
O resumo esperado é \path{e09cfc4793e5394be440906320c7d3008cda5a52b90bbc85bf47446362406af1}. Devem ser registrados data da execução e versões das dependências. O comando de coleta não foi reexecutado para esta revisão textual.

A documentação da revisão bibliográfica está em \path{13-dissertacao/fichamentos/}, no mesmo repositório e referência: o inventário completo identifica as fontes; o manifesto registra versões e localizadores; e o relatório da busca complementar descreve critérios e limites. As tabelas de apoio à caracterização do corpus estão em \path{13-dissertacao/texto-Latex/dados/}. A tabela de cobertura textual e o gráfico de estimativas de segmentação desta dissertação usam, respectivamente, \path{cobertura_textual.csv} e \path{estimativa_chunking.csv}, conferidos com as saídas salvas da análise exploratória.

O artigo que documenta a implementação inicial do protótipo e sua primeira avaliação está em \path{06-inteligencia-artificial-generativa/5-project-paper/chatbot-rag/}, no mesmo repositório. Os artefatos posteriores do experimento de viabilidade — incluindo bateria, configurações e resultados das rodadas — estão em \path{06-inteligencia-artificial-generativa/4-project/}. A Seção~\ref{sec:viabilidade} caracteriza esse percurso como etapa da pesquisa. Seus resultados são preliminares e orientam o protocolo; permanecem distintos dos resultados do experimento principal.
